\chapter{Die Funktion mit acht Argumenten}\label{fall:argumente}

\begin{merksatz}
\boxtitel{Die Vorlage}
\textbf{Zusammenhang.} Auswertung, Python. Formatiert einen Laborwert für die Anzeige. Sechs Aufrufstellen.

\begin{center}
\begin{minipage}{0.92\textwidth}
\code{def formatiere(betrag, exp, einheit, nachkomma, trenner,}\\
\code{\ \ \ \ \ \ \ \ \ \ \ \ \ \ \ sprache, breite, mit\_einheit):}\\
\code{\ \ \ \ ...\ \ \# 34 Zeilen}
\end{minipage}
\end{center}

\begin{center}
\begin{minipage}{0.88\textwidth}
\kursiv{Acht Argumente. Nach Band 20, Kapitel~5 ist die Zahl der Argumente der beste Schnelltest für eine fehlende Zerlegung --- und acht liegt weit über der Schwelle, ab der man hinsehen sollte.}
\end{minipage}
\end{center}
\end{merksatz}

\begin{merksatz}
\boxtitel{Die Diagnose}
\textbf{Wer ruft wie auf?}

\begin{center}
{\small
\renewcommand{\arraystretch}{1.45}
\begin{tabular}{@{}P{2.8cm}P{2.4cm}P{2.6cm}@{}}
\toprule
\textbf{Aufrufstelle} & \textbf{Argumente} & \textbf{Muster}\\
\midrule
Verlaufsansicht & $8$ & \textbf{A}\\
Befundliste & $8$ & \textbf{A}\\
Ausdruck & $8$ & \textcolor{mittel}{B}\\
Ausleitung CSV & $8$ & \textcolor{akzent}{C}\\
Warnungstext & $8$ & A\\
Protokoll & $8$ & \textcolor{akzent}{C}\\
\bottomrule
\end{tabular}}
\end{center}

\begin{center}
\begin{minipage}{0.88\textwidth}
\kursiv{Sechs Aufrufe, drei feste Kombinationen. Das ist der Befund: Die acht Argumente sind nicht acht Freiheitsgrade, sondern drei Fälle, die zu einer Funktion mit Schaltern zusammengezogen wurden --- die verfrühte Zusammenfassung aus Band 20, Kapitel~14.}
\end{minipage}
\end{center}

\textbf{Was die Argumente wirklich sind.}

\begin{center}
\renewcommand{\arraystretch}{1.45}
\begin{tabular}{@{}P{2.8cm}P{2.4cm}P{2.6cm}@{}}
\toprule
\textbf{Gruppe} & \textbf{Argumente} & \textbf{Ist} \\
\midrule
Wert & \textbf{betrag, exp, einheit} & \textbf{eine Größe}\\
Darstellung & nachkomma, trenner, sprache & \textbf{ein Format}\\
Ausgabe & breite, mit\_einheit & Zweck\\
\bottomrule
\end{tabular}
\end{center}

\begin{center}
\kursiv{Drei Gruppen, die immer gemeinsam wandern --- das Erkennungszeichen für einen fehlenden Begriff. Die erste Gruppe ist die \code{Groesse} aus Fall~\ref{fall:einheit}; sie existierte schon und wurde hier auseinandergenommen.}
\end{center}
\end{merksatz}

\begin{technik}
\boxtitel{Das Urteil: überarbeitet}
\begin{center}
\begin{minipage}{0.9\textwidth}
\code{def fuer\_anzeige(wert: Groesse, sprache: Sprache) -> str: ...}\\
\code{def fuer\_druck(wert: Groesse, sprache: Sprache) -> str: ...}\\
\code{def fuer\_ausleitung(wert: Groesse) -> str: ...}
\end{minipage}
\end{center}

\begin{center}
\begin{minipage}{0.86\textwidth}
\kursiv{Drei Funktionen mit ein oder zwei Argumenten statt einer mit acht. Die Formatvorgaben --- Nachkommastellen, Trenner, Breite --- stehen jetzt \kursiv{in} den Funktionen, wo sie hingehören: Sie sind nicht wählbar, sondern folgen aus dem Zweck.}
\end{minipage}
\end{center}
\end{technik}

\begin{merksatz}
\boxtitel{Die Begründung}
\begin{center}
{\small
\renewcommand{\arraystretch}{1.45}
\begin{tabular}{@{}P{2.8cm}P{2.0cm}P{2.2cm}P{2.0cm}@{}}
\toprule
& \textbf{Vorlage} & \textbf{Überarbeitet} & \\
\midrule
Funktionen & $1$ & $\mathbf{3}$ & \\
Zeilen gesamt & $34$ & $\mathbf{41}$ & $+21\,\%$\\
Argumente je Aufruf & $\mathbf{8}$ & $\mathbf{1}$--$\mathbf{2}$ & \\
Zweige & $\mathbf{11}$ & $\mathbf{2}$ & $-82\,\%$\\
Verstehenszeit & $95$ s & $\mathbf{24}$ s & $4{,}0\times$\\
Unzulässige Kombinationen & \textcolor{akzent}{$\mathbf{18}$} & \textcolor{gruen}{$\mathbf{0}$} & \\
\bottomrule
\end{tabular}}
\end{center}

\begin{center}
\begin{minipage}{0.88\textwidth}
\kursiv{Die letzte Zeile ist der eigentliche Gewinn und stand nirgends geschrieben: Von den zweihundertsechsundfünfzig möglichen Argumentkombinationen waren achtzehn sinnlos --- Breite null mit Einheit, CSV-Trenner mit Anzeigesprache. Nach der Trennung gibt es sie nicht mehr.}
\end{minipage}
\end{center}

\textbf{Und die Doppelung?}

\begin{center}
\renewcommand{\arraystretch}{1.45}
\begin{tabular}{@{}P{3.0cm}P{2.2cm}P{2.6cm}@{}}
\toprule
\textbf{Posten} & \textbf{Zeilen} & \\
\midrule
Gemeinsam, ausgelagert & $\mathbf{19}$ & \textcolor{gruen}{Zahl\-formatierung}\\
Je Funktion eigen & $6$--$9$ & \\
\textbf{Tatsächlich doppelt} & $\mathbf{4}$ & \textcolor{gruen}{\textbf{tragbar}}\\
\bottomrule
\end{tabular}
\end{center}

\begin{center}
\kursiv{Vier doppelte Zeilen --- die Rechnung aus Band 22, Kapitel~8: Ein wenig Kopieren ist billiger als eine Fallunterscheidung mit elf Zweigen. Wäre die Doppelung bei dreißig Zeilen gelegen, wäre das Urteil anders ausgefallen.}
\end{center}
\end{merksatz}

\begin{warnung}
\boxtitelr{Der Fehlschluss}
\begin{center}
\renewcommand{\arraystretch}{1.5}
\begin{tabular}{@{}P{3.2cm}P{2.2cm}P{2.4cm}@{}}
\toprule
\textbf{Naheliegend} & \textbf{Richtig?} & \\
\midrule
\enquote{Höchstens drei Argumente} & \textcolor{akzent}{\textbf{nein}} & \textcolor{akzent}{\textbf{Zahl ohne Begründung}}\\
\enquote{Argumente in ein Objekt} & \textcolor{akzent}{\textbf{nein}} & \textbf{verschiebt nur}\\
\enquote{Immer aufteilen} & \textcolor{akzent}{nein} & \textbf{Fall 17}\\
\textbf{\enquote{Gruppen sind fehlende Begriffe}} & \textcolor{gruen}{\textbf{ja}} & \\
\bottomrule
\end{tabular}
\end{center}

\begin{center}
\begin{minipage}{0.88\textwidth}
\kursiv{Die zweite Zeile ist der häufigste Kurzschluss: Acht Argumente in eine Optionsstruktur zu packen macht die Signatur kurz und ändert nichts --- der Aufrufer muss dieselben acht Entscheidungen treffen, und die achtzehn sinnlosen Kombinationen bleiben.}
\end{minipage}
\end{center}

\begin{center}
\kursiv{Der Unterschied: Eine \kursiv{Gruppe}, die immer gemeinsam auftritt und einen Fachnamen hat, wird ein Typ. Eine \kursiv{Sammlung} von Schaltern, die verschiedene Aufrufer verschieden setzen, wird in getrennte Funktionen zerlegt.}
\end{center}
\end{warnung}
\vspace{-4pt}
\noindent Parnas\zit{24} liefert das Kri\-te\-rium, wonach ge\-schnit\-ten wird; Ous\-ter\-hout\zit{23} zeigt an der Zahl der Ar\-gu\-men\-te, wann eine Schnitt\-stel\-le flach ist.
